\documentclass[11pt,a4paper]{article}
\usepackage{fontspec,amsmath,amssymb,amsthm,microtype}
\setmainfont{lmroman10-regular.otf}[BoldFont=lmroman10-bold.otf,ItalicFont=lmroman10-italic.otf,BoldItalicFont=lmroman10-bolditalic.otf]
\usepackage{xcolor}
\usepackage[most]{tcolorbox}
\usepackage{mathrsfs}
\providecommand{\lt}{<}
\providecommand{\gt}{>}
\setmonofont{DejaVuSansMono.ttf}[BoldFont=DejaVuSansMono-Bold.ttf,ItalicFont=DejaVuSansMono-Oblique.ttf,BoldItalicFont=DejaVuSansMono-BoldOblique.ttf,Scale=0.86]
\newfontfamily\leanfallback{STIXTwoMath-Regular.otf}[Scale=0.86]
\definecolor{leanrule}{gray}{0.62}
\definecolor{leanback}{gray}{0.962}
\definecolor{leantitle}{gray}{0.22}
\definecolor{leanmuted}{gray}{0.38}
\newtcolorbox{leanbox}[1]{enhanced,breakable,colback=leanback,colframe=leanrule,
  boxrule=0.4pt,arc=0.6mm,left=2.4mm,right=2.4mm,top=1.8mm,bottom=1.8mm,
  before skip=0.8em,after skip=1em,fonttitle=\small\bfseries,coltitle=leantitle,
  colbacktitle=leanback,titlerule=0.3pt,title={#1}}
\newlength{\leancharwidth}
\newenvironment{leancode}{\par\smallskip\begingroup\ttfamily\footnotesize\raggedright
  \settowidth{\leancharwidth}{M}\setlength{\parskip}{0pt}\setlength{\parindent}{0pt}
  \frenchspacing}{\par\endgroup\smallskip}
\newcommand{\leanline}[2]{\hangindent=\dimexpr#1\leancharwidth+4\leancharwidth\relax
  \hangafter=1\noindent\hspace*{#1\leancharwidth}#2\par}
\newcommand{\leansource}[1]{\par{\footnotesize\color{leanmuted}#1}\par}
\newcommand{\leanlabel}[1]{\par\smallskip\noindent{\small\bfseries #1}\par}
\usepackage[margin=25mm]{geometry}
\usepackage[unicode,hidelinks]{hyperref}
\hypersetup{pdfauthor={LangevinLean Project}}
\setlength{\emergencystretch}{3em}
\setlength{\parskip}{0.4em}
\setlength{\parindent}{0pt}
\allowdisplaybreaks[2]
\raggedbottom
\title{\texorpdfstring{Curvature-width and Gaussian truncation bounds for Hessian differences}{Curvature-width and Gaussian truncation bounds for Hessian differences}}
\author{LangevinLean Project}
\date{}
\begin{document}\maketitle
Draft mathematical exposition; these notes may contain mistakes.


The Hessians of one strongly convex smooth potential differ in operator norm by at most \(M-m\). The proof derives this directly from the centered gradient sector inequality and derivatives. It then proves the identical bound for the random \(L^{2}\) Hessian operators and establishes the precise Gaussian truncation term used in the repaired KLMC2 comparison.



The curvature interval bounds the Hessian difference by its width, while Hessian Lipschitzness gives a displacement bound. Their minimum supports the Gaussian truncation estimate without position–velocity independence.



Earlier results used here: \cite[Section~3.10]{LangevinLean2026ULASC}.

\section{\texorpdfstring{Center the gradient and derive its sharp Lipschitz radius}{Center the gradient and derive its sharp Lipschitz radius}}


For an \(m\)-strongly convex potential with \(M\)-Lipschitz gradient, the strong-smooth interpolation inequality implies \[\left\|\nabla f(x)-\nabla f(y)-\frac{m+M}{2}(x-y)\right\|\le\frac{M-m}{2}\|x-y\|.\] Thus \(g_c(x)=\nabla f(x)-\frac{m+M}{2}x\) is globally \((M-m)/2\)-Lipschitz. The argument constructs the nonnegative Lipschitz constant, rewrites the difference of the centered gradients as the residual in this previously proved sector estimate, and applies that exact estimate. No Hessian or spectral conclusion is inserted as a premise.



Centering at \((m+M)/2\) treats the curvature endpoints symmetrically. The residual Lipschitz radius is half their width. This is sharper than using the full smoothness constant and is especially informative when the potential is nearly quadratic with constant curvature.

\begin{leanbox}{Lean: \texttt{centered\_\allowbreak{}gradient\_\allowbreak{}lipschitz}}
\leansource{\texttt{centered\_\allowbreak{}gradient\_\allowbreak{}lipschitz} \quad KineticLangevin/\allowbreak{}HessianDifference.lean, lines 13–25}
\begin{leancode}
\leanline{0}{theorem\ centered\_gradient\_lipschitz\ (f\ :\ E→ℝ)\ (m\ :\ ℝ)\ (M\ :\ ℝ≥0)}
\leanline{4}{(hf\ :\ Differentiable\ ℝ\ f)\ (hsc\ :\ StrongConvexOn\ univ\ m\ f)}
\leanline{4}{(hLip\ :\ LipschitzWith\ M\ (gradient\ f))\ (hmM\ :\ m≤(M:ℝ))\ :}
\leanline{4}{LipschitzWith\ ⟨((M:ℝ)-m)/2,by\ linarith⟩\ (fun\ x\ ↦\ gradient\ f\ x-((m+(M:ℝ))/2)\ •\ x)}
\end{leancode}
\leanlabel{How the Lean proof follows the argument}
\texttt{centered\_\allowbreak{}gradient\_\allowbreak{}lipschitz} gives the formal counterpart of this step. The main intermediate facts are \texttt{hb}, \texttt{he} in that order. Named proof applications include \texttt{Lipschitz\allowbreak{}With.\allowbreak{}of\_\allowbreak{}dist\_\allowbreak{}le\_\allowbreak{}mul}, \texttt{hb}; the adjacent proof shows their hypotheses and order of use.
\end{leanbox}


Earlier results used here: \cite[Section~3.10.2]{LangevinLean2026ULASC}; \cite[Section~3.10.3]{LangevinLean2026ULASC}.

\section{\texorpdfstring{The Hessian curvature-width bound from derivatives}{The Hessian curvature-width bound from derivatives}}


If \(H_x=D\nabla f(x)\), the derivative of the centered gradient is \(H_x-\frac{m+M}{2}I\). The derivative of a globally Lipschitz map has operator norm at most its Lipschitz constant, so \[\left\|H_x-\frac{m+M}{2}I\right\|\le\frac{M-m}{2}.\] The proof differentiates the centered gradient by subtracting the scalar identity derivative, then applies the derivative/\allowbreak{}Lipschitz operator bound. Subtracting the centered derivatives at two points yields \[\|H_x-H_y\|\le\frac{M-m}{2}+\frac{M-m}{2}=M-m.\] This is the sharp Hessian interval-width estimate. It is derived from the original potential assumptions and gradient derivatives, without an additional continuous-Hessian, symmetry, or operator-range hypothesis.



The derivative of a globally Lipschitz map has operator norm bounded by the same constant wherever it exists. Applying this to the centered gradient bounds each centered Hessian. The triangle inequality between two centered Hessians gives the full interval width \(M-m\).

\begin{leanbox}{Lean: \texttt{hessian\_\allowbreak{}centered\_\allowbreak{}norm\_\allowbreak{}le}, \texttt{hessian\_\allowbreak{}difference\_\allowbreak{}norm\_\allowbreak{}le}}
\leansource{\texttt{hessian\_\allowbreak{}centered\_\allowbreak{}norm\_\allowbreak{}le} \quad KineticLangevin/\allowbreak{}HessianDifference.lean, lines 28–36}
\begin{leancode}
\leanline{0}{theorem\ hessian\_centered\_norm\_le\ (f\ :\ E→ℝ)\ (m\ :\ ℝ)\ (M\ :\ ℝ≥0)}
\leanline{4}{(hf\ :\ Differentiable\ ℝ\ f)\ (hsc\ :\ StrongConvexOn\ univ\ m\ f)}
\leanline{4}{(hLip\ :\ LipschitzWith\ M\ (gradient\ f))\ (hmM\ :\ m≤(M:ℝ))}
\leanline{4}{(x\ :\ E)\ (H\ :\ E→L[ℝ]\ E)\ (hH\ :\ HasFDerivAt\ (gradient\ f)\ H\ x)\ :}
\leanline{4}{‖H-((m+(M:ℝ))/2)\ •\ ContinuousLinearMap.id\ ℝ\ E‖≤((M:ℝ)-m)/2}
\end{leancode}
\leansource{\texttt{hessian\_\allowbreak{}difference\_\allowbreak{}norm\_\allowbreak{}le} \quad KineticLangevin/\allowbreak{}HessianDifference.lean, lines 39–50}
\begin{leancode}
\leanline{0}{theorem\ hessian\_difference\_norm\_le\ (f\ :\ E→ℝ)\ (m\ :\ ℝ)\ (M\ :\ ℝ≥0)}
\leanline{4}{(hf\ :\ Differentiable\ ℝ\ f)\ (hsc\ :\ StrongConvexOn\ univ\ m\ f)}
\leanline{4}{(hLip\ :\ LipschitzWith\ M\ (gradient\ f))\ (hmM\ :\ m≤(M:ℝ))}
\leanline{4}{(x\ y\ :\ E)\ (Hx\ Hy\ :\ E→L[ℝ]\ E)}
\leanline{4}{(hHx\ :\ HasFDerivAt\ (gradient\ f)\ Hx\ x)\ (hHy\ :\ HasFDerivAt\ (gradient\ f)\ Hy\ y)\ :}
\leanline{4}{‖Hx-Hy‖≤(M:ℝ)-m}
\end{leancode}
\leanlabel{How the Lean proof follows the argument}
\texttt{hessian\_\allowbreak{}centered\_\allowbreak{}norm\_\allowbreak{}le} gives the formal counterpart of this step. The main intermediate facts are \texttt{hd} in that order. Named proof applications include \texttt{hd.\allowbreak{}le\_\allowbreak{}of\_\allowbreak{}lipschitz}; the adjacent proof shows their hypotheses and order of use.




\texttt{hessian\_\allowbreak{}difference\_\allowbreak{}norm\_\allowbreak{}le} gives the formal counterpart of this step. The main intermediate facts are \texttt{hx}, \texttt{hy}, \texttt{he} in that order. The final step is the algebraic normalization or ordered arithmetic displayed in the proof.
\end{leanbox}

\section{\texorpdfstring{Identify and bound the random Hessian difference on \(L^{2}\)}{Identify and bound the random Hessian difference on L²}}


The existing random-Hessian construction is a continuous linear map on \(L^{2}\), with representative \(H(X(\omega))Z(\omega)\). One proves that subtracting the two such maps has the exact representative \[(H(X(\omega))-H(Y(\omega)))Z(\omega).\] It rewrites the almost-everywhere representatives of the two random Hessians and their \(L^{2}\) difference. The pointwise sharp bound just proved gives \(\|(H(X)-H(Y))Z\|\le(M-m)\|Z\|\), so the \(L^{2}\) operator norm of the difference is at most \(M-m\). The argument applies the \(L^{2}\) norm monotonicity theorem and the continuous linear operator-norm characterization; no operator bound is assumed on an unspecified random map.



The pointwise difference representation ensures that the operator bound concerns the Hessians frozen at the two random positions. Squaring and integrating the uniform bound transfers it to second mean without any independence assumption between positions and the vector.

\begin{leanbox}{Lean: \texttt{lpHessian\_\allowbreak{}difference\_\allowbreak{}coe}, \texttt{lpHessian\_\allowbreak{}difference\_\allowbreak{}norm\_\allowbreak{}le}}
\leansource{\texttt{lpHessian\_\allowbreak{}difference\_\allowbreak{}coe} \quad KineticLangevin/\allowbreak{}HessianDifference.lean, lines 55–63}
\begin{leancode}
\leanline{0}{theorem\ lpHessian\_difference\_coe\ (g\ :\ E→E)\ (H\ :\ E→E→L[ℝ]\ E)\ (M\ M₂\ :\ ℝ≥0)}
\leanline{4}{(hG\ :\ LipschitzWith\ M\ g)\ (hH\ :\ LipschitzWith\ M₂\ H)}
\leanline{4}{(hg\ :\ ∀\ x,HasFDerivAt\ g\ (H\ x)\ x)\ (X\ Y\ Z\ :\ Lp\ E\ 2\ P)\ :}
\leanline{4}{(lpHessian\ P\ g\ H\ M\ M₂\ hG\ hH\ hg\ X-lpHessian\ P\ g\ H\ M\ M₂\ hG\ hH\ hg\ Y)\ Z\ =ᵐ[P]}
\leanline{6}{fun\ ω\ ↦\ (H\ (X\ ω)-H\ (Y\ ω))\ (Z\ ω)}
\end{leancode}
\leansource{\texttt{lpHessian\_\allowbreak{}difference\_\allowbreak{}norm\_\allowbreak{}le} \quad KineticLangevin/\allowbreak{}HessianDifference.lean, lines 66–77}
\begin{leancode}
\leanline{0}{theorem\ lpHessian\_difference\_norm\_le\ (f\ :\ E→ℝ)\ (H\ :\ E→E→L[ℝ]\ E)\ (m\ :\ ℝ)\ (M\ M₂\ :\ ℝ≥0)}
\leanline{4}{(hf\ :\ Differentiable\ ℝ\ f)\ (hsc\ :\ StrongConvexOn\ univ\ m\ f)}
\leanline{4}{(hG\ :\ LipschitzWith\ M\ (gradient\ f))\ (hmM\ :\ m≤(M:ℝ))}
\leanline{4}{(hH\ :\ LipschitzWith\ M₂\ H)\ (hg\ :\ ∀\ x,HasFDerivAt\ (gradient\ f)\ (H\ x)\ x)\ (X\ Y\ :\ Lp\ E\ 2\ P)\ :}
\leanline{4}{‖lpHessian\ P\ (gradient\ f)\ H\ M\ M₂\ hG\ hH\ hg\ X-lpHessian\ P\ (gradient\ f)\ H\ M\ M₂\ hG\ hH\ hg\ Y‖≤(M:ℝ)-m}
\end{leancode}
\leanlabel{How the Lean proof follows the argument}
\texttt{lp\allowbreak{}Hessian\_\allowbreak{}difference\_\allowbreak{}coe} gives the formal counterpart of this step. Its remaining equalities are discharged by the simplifications displayed in the proof.




\texttt{lp\allowbreak{}Hessian\_\allowbreak{}difference\_\allowbreak{}norm\_\allowbreak{}le} gives the formal counterpart of this step. Named proof applications include \texttt{Continuous\allowbreak{}Linear\allowbreak{}Map.\allowbreak{}op\allowbreak{}Norm\_\allowbreak{}le\_\allowbreak{}bound}, \texttt{Lp.\allowbreak{}norm\_\allowbreak{}le\_\allowbreak{}mul\_\allowbreak{}norm\_\allowbreak{}of\_\allowbreak{}ae\_\allowbreak{}le\_\allowbreak{}mul}; the adjacent proof shows their hypotheses and order of use.
\end{leanbox}

\section{\texorpdfstring{The exact Hessian-difference Gaussian truncation term}{The exact Hessian-difference Gaussian truncation term}}


The original Hessian Lipschitz assumption also gives \(\|H(X)-H(Y)\|\le M_2\|X-Y\|\) pointwise. Combining this with the uniform bound \(M-m\), and retaining a standard Gaussian velocity \(V\), gives \[\|(H(X)-H(Y))V\|_{L^2}\le M_2\sqrt\alpha\,\|X-Y\|_{L^2}+2(M-m)e^{-(\alpha-p)/8},\qquad\alpha\ge5p.\] The argument rewrites the exact random-Hessian difference representative and the position discrepancy representative, combines the two proved operator bounds by their minimum, and applies the Gaussian cutoff theorem. Only the marginal Gaussian law of \(V\) is required; the positions and Hessian difference may depend on \(V\). This is the precise random coefficient bound needed at each free-OU time in the corrected KLMC2 argument.



The global width bound controls the large-velocity region, while Hessian Lipschitzness controls the moderate region through position discrepancy. The Gaussian positive-part tail bounds the excess. Choosing the cutoff later balances absorption of the displacement term against exponential decay of the tail.

\begin{leanbox}{Lean: \texttt{lpHessian\_\allowbreak{}difference\_\allowbreak{}gaussian\_\allowbreak{}truncation}}
\leansource{\texttt{lpHessian\_\allowbreak{}difference\_\allowbreak{}gaussian\_\allowbreak{}truncation} \quad KineticLangevin/\allowbreak{}HessianDifference.lean, lines 82–104}
\begin{leancode}
\leanline{0}{theorem\ lpHessian\_difference\_gaussian\_truncation\ [IsProbabilityMeasure\ P]}
\leanline{4}{(f\ :\ EuclideanSpace\ ℝ\ ι→ℝ)}
\leanline{4}{(H\ :\ EuclideanSpace\ ℝ\ ι→EuclideanSpace\ ℝ\ ι→L[ℝ]\ EuclideanSpace\ ℝ\ ι)}
\leanline{4}{(m\ :\ ℝ)\ (M\ M₂\ :\ ℝ≥0)}
\leanline{4}{(hf\ :\ Differentiable\ ℝ\ f)\ (hsc\ :\ StrongConvexOn\ univ\ m\ f)}
\leanline{4}{(hG\ :\ LipschitzWith\ M\ (gradient\ f))\ (hmM\ :\ m≤(M:ℝ))}
\leanline{4}{(hH\ :\ LipschitzWith\ M₂\ H)\ (hg\ :\ ∀\ x,HasFDerivAt\ (gradient\ f)\ (H\ x)\ x)}
\leanline{4}{(X\ Y\ V\ :\ Lp\ (EuclideanSpace\ ℝ\ ι)\ 2\ P)}
\leanline{4}{(hV\ :\ HasLaw\ V\ (stdGaussian\ (EuclideanSpace\ ℝ\ ι))\ P)}
\leanline{4}{(α\ :\ ℝ)\ (hα\ :\ 5*(Fintype.card\ ι\ :\ ℝ)≤α)\ :}
\leanline{4}{‖(lpHessian\ P\ (gradient\ f)\ H\ M\ M₂\ hG\ hH\ hg\ X-lpHessian\ P\ (gradient\ f)\ H\ M\ M₂\ hG\ hH\ hg\ Y)\ V‖≤}
\leanline{6}{(M₂:ℝ)*Real.sqrt\ α*‖X-Y‖+2*((M:ℝ)-m)*Real.exp\ (-(α-(Fintype.card\ ι\ :\ ℝ))/8)}
\end{leancode}
\leanlabel{How the Lean proof follows the argument}
\texttt{lp\allowbreak{}Hessian\_\allowbreak{}difference\_\allowbreak{}gaussian\_\allowbreak{}truncation} gives the formal counterpart of this step. The main intermediate facts are \texttt{h\allowbreak{}B}, \texttt{h\allowbreak{}L}, \texttt{he} in that order. Named proof applications include \texttt{gaussian\_\allowbreak{}truncation\_\allowbreak{}Lp}; the adjacent proof shows their hypotheses and order of use.
\end{leanbox}

\begin{thebibliography}{ULA-LSI}
\bibitem[ULA-SC]{LangevinLean2026ULASC}
LangevinLean Project. \emph{Unadjusted Langevin under Strong Convexity}. LL-01, 2026. Draft report, content revision d34ead06987c. \url{https://odakyildiz.github.io/LangevinLean/docs/pdf/UnadjustedLangevin/FullProof.pdf}.
\end{thebibliography}

\end{document}
